\chapter{Gauss maps and Hopf differentials}\label{sec:gauss-map-hopf}
Harmonic maps connect CMC surface theory to complex analysis. The sources are \cite{Hel01}, Theorem 2.1, p.~18, and Chapter 4, pp.~34--37. We use the trace $H$ and the sign $h(X,Y)=\langle D_X\nu,Y\rangle$ throughout; the source uses the opposite second fundamental form and averaged mean curvature. Basic complex differentiation and Liouville's theorem are explicit complex-analytic background inputs. Global conformal parametrization of a topological sphere is not proved in this section.

\begin{definition}[Harmonic maps to a unit sphere]\label{def:sphere-valued-harmonic-map}
A smooth map $u:(M,g)\to S^{N-1}\subset\mathbb R^N$ is harmonic if
\[
 \Delta_gu+|du|_g^2u=0.
\]
Equivalently the Euclidean vector $\Delta_gu$ is normal to the target sphere at $u$.
\end{definition}

\begin{proposition}[The sphere-valued energy equation]\label{prop:sphere-energy-equation}
A smooth map $u:(M,g)\to S^{N-1}$ is harmonic precisely when it is stationary for the Dirichlet energy under all smooth target-tangent variations compactly supported in the interior. For each test variation, evaluate $E_U(u)=\frac12\int_U|du|^2$ on a relatively compact domain $U$ containing its support; no finite total energy on a noncompact $M$ is assumed.
\end{proposition}
\begin{proof}
For a field $V$ along $u$ satisfying $\langle V,u\rangle=0$, the variation $(u+tV)/|u+tV|$ stays on the sphere and has initial velocity $V$. Differentiating the energy and integrating by parts gives $E'(0)=-\int\langle\Delta u,V\rangle$. Vanishing for all such fields is exactly the vanishing of the tangential projection of $\Delta u$. Since $\Delta|u|^2=2\langle\Delta u,u\rangle+2|du|^2=0$, its normal component is $-|du|^2u$. This gives the displayed harmonic-map equation. Boundary values are held fixed by the compact support of the test.
\end{proof}

\begin{proposition}[The Gauss-map tension of a Euclidean hypersurface]\label{prop:gauss-map-tension}
For a smooth two-sided hypersurface immersion $F:\Sigma^n\to\mathbb R^{n+1}$ with global unit normal $\nu$, the Gauss map satisfies
\[
 |d\nu|^2=|A|^2,\qquad
 \Delta_\Sigma\nu+|A|^2\nu=\nabla H.
\]
The right side is viewed as a tangent vector in Euclidean space. Thus on a connected hypersurface the Gauss map is harmonic if and only if $H$ is constant. In dimension two this is the Ruh--Vilms statement recalled in \cite{Hel01}, Theorem 2.1.
\end{proposition}
\begin{proof}
The Weingarten formula in our convention is $D_{e_i}\nu=Ae_i$, giving the energy identity. At a point choose a geodesic orthonormal frame. The ambient derivative decomposition gives
\[
 \Delta\nu=\sum_i D_{e_i}(Ae_i)
 =\sum_i(\nabla_{e_i}A)e_i-\sum_i h(e_i,Ae_i)\nu.
\]
The last sum is $-|A|^2\nu$. Contracting Codazzi gives $\sum_i(\nabla_{e_i}A)e_i=\nabla H$. These tensorial formulas hold at every point. The sphere-valued harmonic-map equation therefore reduces to $\nabla H=0$, which is equivalent to constant $H$ on a connected hypersurface.
\end{proof}

\begin{lemma}[Conformal invariance of the harmonic-map equation in two dimensions]\label{lem:surface-harmonic-conformal-invariance}
For a two-dimensional domain, replacing $g$ by $\widetilde g=e^{2\lambda}g$ preserves the harmonicity of a sphere-valued map and its Dirichlet energy. In a conformal coordinate $z=x+iy$ with $g=e^{2\lambda}(dx^2+dy^2)$, the equation is
\[
 u_{xx}+u_{yy}+(|u_x|^2+|u_y|^2)u=0.
\]
\end{lemma}
\begin{proof}
The coordinate divergence expression gives $\Delta_{\widetilde g}=e^{-2\lambda}\Delta_g$ on scalar functions in dimension two, and applies componentwise to $u$. Also $|du|_{\widetilde g}^2=e^{-2\lambda}|du|_g^2$. Thus the equation is multiplied by the same positive factor. The area element is multiplied by $e^{2\lambda}$, canceling the change of energy density. Substitution into a conformal coordinate proves the flat-coordinate form. This cancellation is specific to a two-dimensional domain.
\end{proof}

\begin{proposition}[The Hopf differential of a harmonic map]\label{prop:harmonic-hopf-differential}
For a smooth harmonic map $u$ from a Riemann surface to a unit sphere, the expression
\[
 \Phi_u=\langle u_z,u_z\rangle\,dz^2,
 \qquad \partial_z=\tfrac12(\partial_x-i\partial_y),
\]
is a globally defined holomorphic quadratic differential. The pairing is the complex bilinear extension of the Euclidean inner product, not the Hermitian pairing. It vanishes exactly when $u$ is weakly conformal, meaning $|u_x|=|u_y|$ and $\langle u_x,u_y\rangle=0$ in conformal coordinates. Compare \cite{Hel01}, Chapter 4, p.~34; the source multiplies our coefficient by four.
\end{proposition}
\begin{proof}
For a holomorphic coordinate change $z=z(w)$, the chain rule gives $u_w=u_z z_w$ and hence the coefficient transforms by $z_w^2$, as required for a quadratic differential. By the preceding lemma, $u_{z\bar z}=-(|u_x|^2+|u_y|^2)u/4$. Since $\langle u,u_z\rangle=0$, differentiation gives
\[
 \partial_{\bar z}\langle u_z,u_z\rangle=2\langle u_{z\bar z},u_z\rangle=0.
\]
Finally $4\langle u_z,u_z\rangle=|u_x|^2-|u_y|^2-2i\langle u_x,u_y\rangle$, proving the equivalence with weak conformality. No assertion of immersion follows at points where $du=0$.
\end{proof}

\begin{lemma}[No nonzero holomorphic quadratic differential on the Riemann sphere]\label{lem:sphere-quadratic-differential-zero}
Every holomorphic quadratic differential on the standard Riemann sphere vanishes identically. Therefore every smooth harmonic map from that sphere to a unit sphere is weakly conformal.
\end{lemma}
\begin{proof}
Write the differential in the finite chart as $q(z)dz^2$, where $q$ is entire. In the coordinate $w=1/z$ at infinity its coefficient is $\widetilde q(w)=q(1/w)w^{-4}$. Smoothness and holomorphicity near $w=0$ make this coefficient bounded there. Hence $q(z)=z^{-4}\widetilde q(1/z)=O(|z|^{-4})$ at infinity. Liouville's theorem implies that the bounded entire function $q$ is constant, and the decay makes the constant zero. Apply Proposition \ref{prop:harmonic-hopf-differential} to obtain the last assertion. This is the argument of \cite{Hel01}, Theorem 4.1, p.~35.
\end{proof}

\begin{theorem}[Hopf rigidity for a conformally parametrized CMC sphere]\label{thm:hopf-cmc-conformal-sphere}
Let $F:S^2\to\mathbb R^3$ be a smooth conformal immersion from the standard Riemann sphere, with constant scalar mean curvature $H$. Then $H\ne0$ and its image is a round sphere of radius $2/|H|$. This is \cite{Hel01}, Theorem 4.2, pp.~36--37, with the conformal-parametrization hypothesis stated explicitly.
\end{theorem}
\begin{proof}
The global orientation gives a smooth Gauss map. By Proposition \ref{prop:gauss-map-tension} and Lemma \ref{lem:surface-harmonic-conformal-invariance}, it is harmonic for the standard conformal structure and hence weakly conformal by Lemma \ref{lem:sphere-quadratic-differential-zero}. At a principal orthonormal frame, $d\nu=A$, so weak conformality says $\kappa_1^2=\kappa_2^2$. If $H=0$, the identity $\Delta F=\vec H=0$ makes every coordinate function harmonic on a compact connected surface and therefore constant, contradicting immersion. Thus $H\ne0$. The equality of squares together with $\kappa_1+\kappa_2=H\ne0$ implies $\kappa_1=\kappa_2=H/2$ at every point.

Put $\lambda=H/2$. Since $\lambda$ is constant and nonzero, differentiating $F-\nu/\lambda$ gives zero. Its value is a fixed center $c$, and $|F-c|=1/|\lambda|=2/|H|$. The image is open in that sphere by the inverse function theorem and closed by compactness, so it is the entire sphere. This proves the claimed image statement without assuming embeddedness.
\end{proof}

\begin{remark}[The topological-sphere formulation and its analytic input]\label{rem:hopf-uniformization-input}
To conclude the same theorem for a smooth CMC immersion of an abstract surface merely known to be diffeomorphic to $S^2$, apply the genus-zero uniformization theorem to its induced metric and reparametrize conformally. Uniformization is an external analytic input here; a topological identification alone is not a conformal parametrization. The local geometric proof above remains valid for every immersed conformal sphere and does not classify CMC immersions of positive genus.
\end{remark}
